\documentclass[../differential_geometry.tex]{subfiles}
\begin{document}
\section{Aula 04 - de Fevereiro, 2026}
\subsection{Motivações}
\begin{itemize}
	\item Ponto de Inflexão e Vértice;
	\item Teorema Fundamental de Curvas Planas;
	\item Forma Canônica;
	\item Contato entre Curvas;
	\item Contato de uma Curva com Retas.
\end{itemize}
\subsection{Teorema Fundamental de Curvas Planas}
Na aula passada, terminamos definindo vértices, pontos de inflexão e estudando alguns exemplos; hoje, daremos continuidade a esse tema, ilustrando mais alguns exemplos e chegando ao importante \hyperlink{fundamental_plane_curve}{\textit{Teorema Fundamental da Teoria Local de Curvas Planas}}.
\begin{example}[Elipse]
	A elipse pode ser parametrizada por
	\[
		\alpha (u)=(a\cos^{}{u}, b\sin^{}{u}), \quad u\in \mathbb{R},\; a, b>0,\; a\neq b.
	\]
	Como
	\begin{align*}
		 & \alpha '(u)=(-a \sin^{}{u}, b \cos^{}{u})  \\
		 & \alpha ''(u)=(-a\cos^{}{u}, -b\sin^{}{u}),
	\end{align*}
	a curvatura da elipse em um ponto u é dada pela relação
	\[
		\kappa (u)=\frac{ab}{(a^{2}\sin^{2}{u}+b^{2}\cos^{2}{u})^{\frac{3}{2}}}.
	\]
	Observe aqui que \(\kappa (u)>0\) para todo u, ou seja, a \textbf{elipse não possui pontos de inflexão!}.

	Por outro lado, derivando a curva, obtemos
	\[
		\kappa '(u)= - \frac{3ab(a^{2}-b^{2})\sin^{}{u}\cos^{}{u}}{(a^{2}\sin^{2}{u}+b^{2}\cos^{2}{u})^{\frac{5}{2}}};
	\]
	daí, os vértices da elipse podem ser encontrados por
	\[
		\kappa '(u)=0 \Longleftrightarrow u=0,\; \frac{\pi }{2},\; \pi ,\; \frac{3\pi }{2}.
	\]
	Portanto, a elipse possui quatro vértices.
\end{example}
\begin{example}[Parábola]
	A parábola pode ser parametrizada por \(\alpha (u)=(u, u^{2}),\; u\in \mathbb{R}\), da qual obtemos
	\begin{align*}
		 & \tikz[baseline=-0.5ex]\draw[black, fill=black, radius=1.5pt](0, 0)circle; \; \alpha '(u)=(1,2u)   \\
		 & \tikz[baseline=-0.5ex]\draw[black, fill=black, radius=1.5pt](0, 0)circle; \; \alpha''(u) = (0,2),
	\end{align*}
	donde concluímos que
	\[
		\kappa (u) = \frac{2}{(1+4u^{2})^{\frac{3}{2}}},
	\]
	que, novamente, nunca se anula -- a parábola também não tem pontos de inflexão! No entanto, novamente,
	\[
		\kappa '(u) = - \frac{24u}{(1+4u^{2})^{\frac{5}{2}}} \Rightarrow \kappa '(u)=0 \Longleftrightarrow u =0.
	\]
	Portanto, a parábola tem um único ponto de vértice em \(\alpha (0).\)
\end{example}


\end{document}
